% RQ1(b): Baseline failure, full-controller intervention timing,
% short-burst usability, and sustained behavior.
%
% Results are aggregated from the seven 2026-07-27 CaptureMetrics workbooks
% according to the protocol in the implementation repository.
% Implementation reference inspected at commit:
% 267a29c769428df16c7b96bd18956764da7c37db
%
% Timeout onset (E/M) repeats the no-controller first-timeout index from
% Table~\ref{tab:timeout_index} (M+S columns) for the two endpoint scenarios,
% matching the Baseline arm of Table~\ref{tab:rq1_ablation}.  It is a
% controller-off reference, not part of the proposed implementation, and is
% separated from every full-controller column by its own top-level header and
% a double vertical rule.  E and M each get their own column so that the shot
% indices line up down the column; every action cell reports its absolute onset
% followed by [Delta] relative to the Timeout onset column.  Negative and positive
% shifts mean earlier and later controller action, respectively; they do not
% encode better and worse behavior because admission and pacing coordinate.
% The two timeout columns are named as one pair: the controller-off arm reports
% Timeout onset (when it first failed) and the full controller reports Timeout
% count (how often it failed at all).
% The full controller incurred no Capture Timeout within the 30-shot horizon
% in any row.  This repeats the Ours column of Table~\ref{tab:rq1_ablation},
% where the same result appears as an unreached first-timeout statistic; the
% Timeout count column restates it per row so that Deadline safety carries the
% outcome next to its margin and this table can be read on its own.  Slack
% P5 is normalized by the product Capture Timeout deadline.  The deadline
% constant itself is internal and must not appear in the manuscript, so only
% the normalized percentage is reported.
% Deadline safety sits immediately right of the controller-off reference so
% that the failure the controller has to prevent and the outcome it achieved
% read across in one line; intervention onset, work completion, and pacing
% cost then explain how that outcome was reached and what it cost.
% Draft completion and pacing cost are reported at @5 and @30.  The @5 prefix
% characterizes short-burst behavior close to the user-facing capture
% experience, whereas @30 characterizes sustained behavior.

\begin{table*}[t]
  \centering
  \setlength{\abovecaptionskip}{2pt}%
  \setlength{\belowcaptionskip}{3pt}%
  % The device is now a column, so it no longer belongs in the caption.
  \caption{RQ1(b) controller-off failure reference and full-controller
  behavior.}
  \label{tab:rq1_preventive_alignment}
  \label{tab:rq1_controller_behavior}

  {\scriptsize
  % One value per cell.  E and M have their own columns and the shift sits to
  % the right of the absolute onset; both parts are set in fixed-width boxes
  % so that digits and brackets line up down the column.  @5 and @30 are
  % likewise separate columns, so no in-cell separator is needed anywhere.
  \newlength{\rqOnsetAbsW}\settowidth{\rqOnsetAbsW}{30}%
  \newlength{\rqOnsetDelW}\settowidth{\rqOnsetDelW}{[+6]}%
  \newcommand{\rqOnset}[2]{%
    \makebox[\rqOnsetAbsW][r]{#1}\,\makebox[\rqOnsetDelW][r]{[#2]}}%
  % Not reached within the horizon.  Holding the dash in the absolute slot and
  % padding an empty bracket slot after it left the dash stranded mid-cell in a
  % right-aligned column, so it is set flush right across the whole onset box
  % instead.  The width is measured from the widest real cell so it tracks any
  % change to the two slots above.
  \newlength{\rqOnsetW}\settowidth{\rqOnsetW}{\rqOnset{30}{+6}}%
  \newcommand{\rqNoOnset}{\makebox[\rqOnsetW][r]{--}}%
  % Controller-off E and M hold two characters, but their merged header is far
  % wider.  A \multicolumn wider than the columns it spans dumps the whole
  % excess into the last spanned column, which would leave E cramped against
  % the left rule and M adrift.  Sizing both cells to half the header width
  % removes the excess, so each label sits centred over its own column while
  % the digits stay right aligned inside the box.
  \newlength{\rqShotW}\settowidth{\rqShotW}{\textbf{Timeout onset}}%
  \setlength{\rqShotW}{0.5\rqShotW}%
  \newcommand{\rqShot}[1]{\makebox[\rqShotW][r]{#1}}%
  \newcommand{\rqShotHd}[1]{\makebox[\rqShotW][c]{#1}}%
  \setlength{\arrayrulewidth}{0.35pt}%
  \setlength{\doublerulesep}{0.8pt}%
  \renewcommand{\arraystretch}{1.08}%
  % \fittabcolsep (macros.tex) solves for the \tabcolsep that makes the outer
  % rules land exactly on \textwidth, so the right border meets the text
  % margin without an eyeballed constant.
  \newsavebox{\rqtabbox}%
  \newcommand{\rqtabular}{%
  \begin{tabular}{|c|c|c||cc||r|r||rr|rr||rr|rr||rr|rr|}
    \hline\hline
    % First header row: the two arms.  The controller-off arm has no metric
    % family of its own, so its arm name spans rows 1--2 instead of leaving a
    % blank band or repeating itself.  That drops its metric name onto row 3
    % with every other metric name, which is what keeps Timeout onset and
    % Timeout count reading as one pair on the same line.
    \multirow{4}{*}[-1.8ex]{\textbf{Device}} &
    \multirow{4}{*}[-1.8ex]{\textbf{Condition}} &
    \multirow{4}{*}[-1.8ex]{\makecell[c]{\textbf{Starting}\\\textbf{overheat}\\\textbf{level}}} &
    \multicolumn{2}{c||}{\multirow{2}{*}{\makecell[c]{%
      \textbf{Controller-off}\\\textbf{reference}}}} &
    \multicolumn{14}{c|}{\textbf{Our-controller behavior}} \\
    \cline{6-19}
    % Second header row: metric families.  Columns 4--5 stay covered above.
    & & &
    \multicolumn{2}{c||}{} &
    \multicolumn{2}{c||}{\textbf{Deadline safety}} &
    \multicolumn{4}{c||}{\textbf{Intervention onset}} &
    \multicolumn{4}{c||}{\textbf{Draft work completion}} &
    \multicolumn{4}{c|}{\textbf{Pacing cost}} \\
    \cline{4-19}
    % Third header row: one metric name per family, all on this line.  The
    % @5/@30 and E/M splits are pushed to the fourth row.  Timeout count and
    % Slack P5 have no such split, so the cline below skips them and their
    % cells run on into row 4; they are \multirow so the label is centred in
    % that taller box rather than clinging to its top edge.  Row 3 is two
    % lines tall and row 4 one, so multirow's equal-height assumption is off
    % by half a line and the vmove below corrects it.
    & & &
    \multicolumn{2}{c||}{\textbf{Timeout onset}} &
    % Deadline safety carries two metrics: the outcome and the margin.  The
    % Timeout count column restates the RQ1(a) result per row so that this
    % table is readable on its own; it is zero throughout by construction.
    % It is left bare so that it matches Timeout onset; the 30-shot horizon
    % it counts over is stated in the note below the table.
    \multicolumn{1}{c|}{\multirow{2}{*}[-0.4ex]{%
      \makecell[c]{\textbf{Timeout}\\\textbf{count}}}} &
    \multicolumn{1}{c||}{\multirow{2}{*}[-0.4ex]{%
      \makecell[c]{\textbf{Slack P5}\\\textbf{(\%)}}}} &
    \multicolumn{2}{c|}{\textbf{Admission-skip}} &
    \multicolumn{2}{c||}{\textbf{Pacing-delay}} &
    % \textbf does not reach math mode, so bold these through \boldmath to
    % match the surrounding header weight.
    \multicolumn{2}{c|}{\makecell[c]{\mbox{\boldmath$M{+}S$}\\\textbf{(\%)}}} &
    \multicolumn{2}{c||}{\makecell[c]{\mbox{\boldmath$M$}\\\textbf{(\%)}}} &
    \multicolumn{2}{c|}{\makecell[c]{\textbf{Activated}\\\textbf{(\%)}}} &
    \multicolumn{2}{c|}{\makecell[c]{\textbf{Delay P50}\\\textbf{(ms)}}} \\
    \cline{4-5}\cline{8-19}
    % Fourth header row: E/M under the shot-index families, @5/@30 under the
    % completion and pacing-cost families.  All 19 columns must be present or
    % the row shifts left; columns 6 and 7 are intentionally empty.
    & & &
    \multicolumn{1}{c}{\rqShotHd{\textbf{E}}} &
    \multicolumn{1}{c||}{\rqShotHd{\textbf{M}}} &
    & &
    \multicolumn{1}{c}{\textbf{E}}  & \multicolumn{1}{c|}{\textbf{M}} &
    \multicolumn{1}{c}{\textbf{E}}  & \multicolumn{1}{c||}{\textbf{M}} &
    \multicolumn{1}{c}{\textbf{@5}} & \multicolumn{1}{c|}{@30} &
    \multicolumn{1}{c}{\textbf{@5}} & \multicolumn{1}{c||}{@30} &
    \multicolumn{1}{c}{\textbf{@5}} & \multicolumn{1}{c|}{@30} &
    \multicolumn{1}{c}{\textbf{@5}} & \multicolumn{1}{c|}{@30} \\
    \hline\hline
    % One device is measured, so the Device cell spans all fourteen data rows;
    % the block separator below is \cline{2-19} so it does not cut through it.
    \multirow{14}{*}{\makecell[c]{S26\\Ultra}}
      & \multirow{7}{*}{\makecell[c]{12MP\\normal}}
      & Lv0 & \rqShot{--} & \rqShot{--} & 0 & 30.0
        & \rqNoOnset & \rqNoOnset & \rqNoOnset & \rqNoOnset
        & \textbf{100} & 100 & \textbf{100} & 100
        & \textbf{0} & 0 & \textbf{--} & -- \\
    \cline{3-19}
      & & Lv1 & \rqShot{--} & \rqShot{--} & 0 & 21.9
        & \rqNoOnset & \rqNoOnset & \rqNoOnset & \rqNoOnset
        & \textbf{100} & 100 & \textbf{100} & 100
        & \textbf{0} & 0 & \textbf{--} & -- \\
    \cline{3-19}
      & & Lv2 & \rqShot{24} & \rqShot{27} & 0 & 4.7
        & \rqOnset{23}{-1} & \rqOnset{29}{+2}
        & \rqOnset{27}{+3} & \rqOnset{28}{+1}
        & \textbf{100} & 88.9 & \textbf{100} & 90.0
        & \textbf{0} & 9.2 & \textbf{--} & 182 \\
    \cline{3-19}
      & & Lv3 & \rqShot{13} & \rqShot{18} & 0 & 2.9
        & \rqOnset{9}{-4}  & \rqOnset{18}{0}
        & \rqOnset{12}{-1} & \rqOnset{15}{-3}
        & \textbf{100} & 57.8 & \textbf{100} & 61.5
        & \textbf{0} & 30.3 & \textbf{--} & 372 \\
    \cline{3-19}
      & & Lv4 & \rqShot{10} & \rqShot{17} & 0 & 5.1
        & \rqOnset{4}{-6}  & \rqOnset{15}{-2}
        & \rqOnset{10}{0}  & \rqOnset{14}{-3}
        & \textbf{96.0} & 47.0 & \textbf{96.0} & 47.0
        & \textbf{0} & 33.4 & \textbf{--} & 386 \\
    \cline{3-19}
      & & Lv5 & \rqShot{8} & \rqShot{10} & 0 & 3.8
        & \rqOnset{9}{+1}  & \rqOnset{15}{+5}
        & \rqOnset{8}{0}   & \rqOnset{10}{0}
        & \textbf{100} & 55.6 & \textbf{100} & 55.6
        & \textbf{0} & 44.8 & \textbf{--} & 451 \\
    \cline{3-19}
      & & Lv6 & \rqShot{7} & \rqShot{10} & 0 & 2.7
        & \rqOnset{9}{+2}  & \rqOnset{16}{+6}
        & \rqOnset{8}{+1}  & \rqOnset{9}{-1}
        & \textbf{100} & 47.0 & \textbf{100} & 48.2
        & \textbf{0} & 41.7 & \textbf{--} & 519 \\
    \cline{2-19}
      & \multirow{7}{*}{\makecell[c]{24MP\\memory\\pressure}}
      & Lv0 & \rqShot{26} & \rqShot{30} & 0 & 3.7
        & \rqOnset{26}{0}  & \rqOnset{30}{0}
        & \rqOnset{25}{-1} & \rqOnset{26}{-4}
        & \textbf{100} & 93.3 & \textbf{100} & 93.3
        & \textbf{0} & 11.5 & \textbf{--} & 243 \\
    \cline{3-19}
      & & Lv1 & \rqShot{19} & \rqShot{22} & 0 & 5.1
        & \rqOnset{20}{+1} & \rqOnset{20}{-2}
        & \rqOnset{22}{+3} & \rqOnset{22}{0}
        & \textbf{100} & 81.7 & \textbf{100} & 81.7
        & \textbf{0} & 24.1 & \textbf{--} & 310 \\
    \cline{3-19}
      & & Lv2 & \rqShot{14} & \rqShot{21} & 0 & 4.4
        & \rqOnset{15}{+1} & \rqOnset{18}{-3}
        & \rqOnset{15}{+1} & \rqOnset{18}{-3}
        & \textbf{100} & 57.8 & \textbf{100} & 58.9
        & \textbf{0} & 29.9 & \textbf{--} & 269 \\
    \cline{3-19}
      & & Lv3 & \rqShot{7} & \rqShot{10} & 0 & 6.1
        & \rqOnset{7}{0}   & \rqOnset{10}{0}
        & \rqOnset{7}{0}   & \rqOnset{11}{+1}
        & \textbf{100} & 31.7 & \textbf{100} & 31.7
        & \textbf{0} & 19.2 & \textbf{--} & 446 \\
    \cline{3-19}
      & & Lv4 & \rqShot{5} & \rqShot{9} & 0 & 6.3
        & \rqOnset{3}{-2}  & \rqOnset{12}{+3}
        & \rqOnset{7}{+2}  & \rqOnset{8}{-1}
        & \textbf{88.0} & 45.0 & \textbf{88.0} & 45.3
        & \textbf{0} & 18.3 & \textbf{--} & 259 \\
    \cline{3-19}
      & & Lv5 & \rqShot{4} & \rqShot{8} & 0 & 8.2
        & \rqOnset{5}{+1}  & \rqOnset{10}{+2}
        & \rqOnset{7}{+3}  & \rqOnset{11}{+3}
        & \textbf{96.0} & 32.0 & \textbf{96.0} & 32.0
        & \textbf{0} & 19.3 & \textbf{--} & 396 \\
    \cline{3-19}
      & & Lv6 & \rqShot{4} & \rqShot{8} & 0 & 7.0
        & \rqOnset{4}{0}   & \rqOnset{10}{+2}
        & \rqOnset{7}{+3}  & \rqOnset{11}{+3}
        & \textbf{90.0} & 40.8 & \textbf{90.0} & 40.8
        & \textbf{0} & 27.6 & \textbf{--} & 369 \\
    \hline\hline
  \end{tabular}}%
  \fittabcolsep{\rqtabbox}{\rqtabular}{\textwidth}{38}%
  \rqtabular
  }

  \par\vspace{2pt}
  {\scriptsize Bold @5 values characterize the first five shots as a
  short-burst usability proxy; @30 values characterize the full burst.
  The controller-off Timeout onset columns are references; all columns to their
  right report the full controller.  Timeout count reports Capture Timeouts
  across the included full-controller runs within the 30-shot horizon, and
  Slack P5 reports the remaining lower-tail deadline margin.  Onset cells
  report the absolute shot index followed by its bracketed shift from the
  corresponding controller-off Timeout onset.  Negative and positive
  shifts denote earlier and later action, respectively.
  -- indicates that the event was not observed within the corresponding
  horizon.  Pacing delay is the median over positive applied delays.}
\end{table*}
